%% OVERFORE-Subdocumento1.tex — generado desde OVERFORE-Subdocumento1.docx
%% Formato: clase overfore, modo subdoc. El formato lo pone la clase; aquí sólo va el contenido.
\documentclass[subdoc]{../../comun/overfore}
\numsubdoc{1}
\nombresubdoc{Presentación de la empresa}
\informe{Informe Preparatorio N° 2}
\fecha{12 de octubre de 2026}

\begin{document}

\capitulo{Introducción}
Overfore S.A. postula a la Licitación N° TFEP-01/2026 para implantar y operar SITCO Aranda en la faena de la Compañía Minera Altos de Aranda S.A. Esta oferta recoge la trayectoria de la empresa, organización y gobierno interno, los proyectos ejecutados, las certificaciones y las alianzas tecnológicas.

Los cuatro proyectos de referencia se detallan en el Formulario T-6, con las cartas de sus mandantes en los Anexos A a D del Subdocumento 1. El equipo clave se detalla en el Anexo E del Subdocumento 1, y la dotación de apoyo y los especialistas externos, en el Anexo F del Subdocumento 1.

\seccion{Presentación de la empresa}
Overfore S.A. es una consultora de desarrollo de sistemas que diseña, implanta y opera plataformas digitales de misión crítica para operaciones industriales. Se constituyó en 2018, tiene su casa matriz en Santiago y una base de soporte en Antofagasta para la zona norte. La Tabla \ref{tab:datos-institucionales-overfore-s} reúne sus datos institucionales.

\begin{tabla}{I{0.396}I{0.604}}{Datos institucionales de Overfore S.A.}{tab:datos-institucionales-overfore-s}
  {\cab{Campo} & \cab{Dato institucional}}
  {Fuente: elaboración propia.}
Razón social & Overfore S.A. \\
RUT & 77.482.961-K \\
Año de constitución & 2018 \\
Representante legal & Claudio Cabello Vásquez \\
Casa matriz & Santiago de Chile \\
Cobertura & Chile, soporte de la zona norte desde Antofagasta \\
Licitaciones y proyecto Aranda & licitaciones@overfore.cl \\
\end{tabla}
Claudio Cabello Vásquez representa legalmente a Overfore, firma esta oferta y dirige el proyecto. La base de Antofagasta, a 145 km de la faena de Aranda (Compañía Minera Altos de Aranda [CMAA], 2026b, cap. 15, p. 26), ya presta soporte presencial en la zona norte.

La trayectoria de la empresa se resume en cuatro plataformas ejecutadas entre 2021 y 2025 como contratista principal, en cadena de frío, manufactura de proceso, logística portuaria y minería del cobre. Las cuatro siguen en operación a cargo de su centro de operaciones, con mesa de servicio, monitoreo continuo y soporte de segundo nivel.

Los servicios se ordenan en ocho líneas de negocio que cubren el ciclo completo de una plataforma: arquitectura híbrida, integración industrial segura, trazabilidad y conciliación, captura en terreno sin cobertura, datos y analítica operacional, seguridad y resiliencia, infraestructura de borde y terreno, e implantación y operación. En cada línea la empresa presta tres servicios: diseño e integración de la solución, implantación en terreno y operación continua con nivel de servicio comprometido. La solvencia de la empresa se acredita con sus estados financieros auditados de los últimos tres ejercicios, el certificado bancario de capacidad de emisión de garantías y el Boletín Laboral y Previsional, que se entregan en el Sobre N° 1.

Las capacidades instaladas se agrupan en seis disciplinas. La Tabla \ref{tab:capacidades-instaladas-respaldo} indica qué cubre cada una y qué la respalda: un proyecto en operación, un responsable nominado o una alianza.

\begin{tabla}{I{0.233}I{0.384}I{0.384}}{Capacidades instaladas y su respaldo.}{tab:capacidades-instaladas-respaldo}
  {\cab{Disciplina} & \cab{Qué cubre} & \cab{Respaldo}}
  {Fuente: elaboración propia.}
Arquitectura híbrida y continuidad & Componentes en nube y en faena, operación degradada y sincronización diferida & RedFrío 360 y FaenaLink, Arquitecto de Solución, NubeAndes \\
Integración industrial segura & ERP, despacho, laboratorio, telemetría e historiador en solo lectura & Planta Conecta, Líder de Integración, Integra OT \\
Infraestructura de borde y terreno & Dispositivos, lectores, gabinetes y servidores industriales & RedFrío 360 y FaenaLink, base de Antofagasta, RedNorte \\
Seguridad y resiliencia & Segmentación, acceso, cifrado, respaldos, alta disponibilidad y recuperación & ISO/IEC 27001 e ISO 22301, Encargado de Seguridad de la Información \\
Datos, trazabilidad y analítica & Trazabilidad, conciliación, auditoría e indicadores por turno & NexoPuerto y FaenaLink, Líder de Datos \\
Implantación y operación & Ambientes, pruebas, marcha blanca, transferencia y soporte en faena & Cuatro plataformas en operación, ISO/IEC 20000-1 \\
\end{tabla}
Cada disciplina tiene al menos un proyecto en operación y un responsable o socio que la sostiene en Aranda.

\seccion{Estructura Organizacional}
La empresa se organiza en una dirección, dos unidades de control independientes y tres áreas de línea. La Figura \ref{fig:estructura-organizacional-overfore-s} muestra esa estructura y lo que cada área asigna al proyecto.

\figura{15.74}{fig-1-1-estructura-organizacional-overfore-s}{Estructura organizacional de Overfore S.A.}{fig:estructura-organizacional-overfore-s}{Fuente: elaboración propia.}
La dirección recae en el representante legal, que también dirige el proyecto. La oficina de calidad y el responsable de seguridad de la información dependen de la dirección y no del área de desarrollo. Así, quien construye no aprueba su propio trabajo.

Las tres áreas de línea son el origen de la dotación de apoyo. El área de desarrollo y el centro de operaciones están en la casa matriz, y la base de la zona norte, en Antofagasta. El centro de operaciones y la base ya sostienen las cuatro plataformas en operación, de modo que la operación de Aranda amplía un servicio que la empresa ya presta.

Overfore cuenta con 64 personas, 28 en el área de desarrollo, 26 en el centro de operaciones y 10 en la base de Antofagasta. El centro de operaciones asigna a Aranda 15 analistas, diez de mesa y NOC y cinco del SOC. Seis provienen de su dotación actual y nueve se contratan antes del mes 21, de modo que las cuatro plataformas en operación conservan 20 analistas.

La empresa asigna a Aranda nueve integrantes del equipo clave y personal propio de apoyo. Ese apoyo llega a treinta personas en la implementación, en los meses 14 y 15, y a veintiuna en la operación, con veintiséis entre los meses 21 y 23. A ellos se agregan doce perfiles de especialistas externos contratados para trabajos puntuales. Con el equipo clave y los especialistas externos de ese mes, el proyecto llega a 45 personas distintas en el mes 14, el máximo del contrato (Formulario T-15). El detalle por perfil, cantidad y período está en el Anexo F del Subdocumento 1.

\seccion{Gobierno interno Calidad, Seguridad y Conocimiento}
El gobierno interno distingue cinco ámbitos, cada uno con su propio control: calidad, seguridad de la información, configuración y cambios, gestión del conocimiento, y continuidad y operación. La Tabla \ref{tab:ambitos-gobierno-interno} indica, para cada uno, su política, la instancia que la revisa y su responsable en el proyecto.

\begin{tabla}{I{0.200}I{0.366}I{0.222}I{0.212}}{Ámbitos del gobierno interno.}{tab:ambitos-gobierno-interno}
  {\cab{Ámbito} & \cab{Política} & \cab{Instancia} & \cab{Responsable en Aranda}}
  {Fuente: elaboración propia.}
Calidad & Revisión de entregables por hito, métricas de defectos y acciones correctivas (ISO 9001:2015) & Oficina de calidad, independiente del desarrollo & Líder de Calidad y Pruebas \\
Seguridad de la información & Sistema de gestión certificado en ISO/IEC 27001:2022, con revisiones de acceso y de vulnerabilidades & Responsable de seguridad independiente & Encargado de Seguridad de la Información \\
Configuración y cambios & Repositorios versionados, revisión por pares y ambientes separados & Comité de Arquitectura & Líder de Desarrollo \\
Gestión del conocimiento & Base documental controlada y manuales en español, actualizados al cierre de cada hito & Comité de Proyecto & Jefe de Proyecto \\
Continuidad y operación & Monitoreo, respaldos probados y ejercicios de recuperación (ISO 22301:2019) & Comité de Operación & Líder de Operación / SRE \\
\end{tabla}
La calidad y la seguridad no quedan a cargo de quien ejecuta el trabajo que controlan. Cada ámbito tiene además una instancia que lo revisa con regularidad y un responsable nominado en el proyecto, que responde por su cumplimiento ante el cliente.

\seccion{Experiencia y Certificaciones}
La experiencia se acredita con cuatro proyectos de software de misión crítica, terminados y en operación, ejecutados en los últimos cinco años como contratista principal. El detalle de cada uno está en el Formulario T-6, y la carta de su mandante, en los Anexos A a D del Subdocumento 1. La Tabla \ref{tab:proyectos-referencia} los resume.

\begin{tabla}{I{0.223}I{0.163}I{0.112}I{0.196}I{0.306}}{Proyectos de referencia.}{tab:proyectos-referencia}
  {\cab{Proyecto y mandante} & \cab{Industria} & \cab{Período} & \cab{Arquitectura y servicio} & \cab{Volumen}}
  {Fuente: elaboración propia.}
FaenaLink · Minera Pampa Tarel SpA & Minería del cobre & 2024 a 2025 & Híbrida de 99,5 \% mensual & 46 camiones, 3.200 personas acreditadas y 1,4 millones de eventos al año \\
NexoPuerto · Terminal Logístico Bahía Sur S.A. & Logística portuaria & 2023 a 2024 & Híbrida de 99,7 \% mensual & 1,8 millones de eventos al año, 170 usuarios y 42 equipos móviles \\
Planta Conecta · Alimentos Cordillera SpA & Manufactura de proceso & 2022 a 2023 & Híbrida de 99,6 \% mensual & 96 activos, 3.500 etiquetas de proceso y 520.000 eventos al año \\
RedFrío 360 · Frío Austral Operaciones Ltda. & Cadena de frío & 2021 a 2022 & Híbrida de 99,5 \% mensual & 24 centros, 680 dispositivos y 2,4 millones de lecturas al año \\
\end{tabla}
Los cuatro proyectos están terminados y en operación, son híbridos y mantienen una disponibilidad de 99,5 \% o más.

FaenaLink es la referencia en la minería del cobre, la industria del cliente. Se ejecutó en una faena de cobre en operación, con turnos 7x7, acreditación de contratistas y registro en terreno sin cobertura. El personal de la empresa trabajó acreditado en la faena, bajo el Reglamento de Seguridad Minera, y la base de Antofagasta mantiene hoy su soporte presencial. Su escala es comparable a la de Aranda, que opera 38 camiones y tiene 4.180 personas con acceso a la faena en régimen, frente a los 46 camiones y las 3.200 personas acreditadas de FaenaLink (CMAA, 2026b, cap. 2, p. 5, y cap. 14, p. 23). En monto, en cambio, los cuatro proyectos son menores que Aranda: el mayor, FaenaLink, se ejecutó por entre USD 650.000 y USD 800.000 (Formulario T-6). Lo que respalda la oferta no es el tamaño de un contrato anterior, sino que la empresa ya construyó y opera cada parte del alcance de Aranda: la arquitectura híbrida con operación sin enlace (FaenaLink y RedFrío 360), la lectura de la red de control sin escribir en ella (Planta Conecta), la trazabilidad por unidad (NexoPuerto) y el soporte presencial en faena con turnos 7x7 (FaenaLink). Aranda reúne esas partes en un solo contrato de 56 meses, por lo que la empresa amplía su dotación con contrataciones programadas (Anexo F) y suma socios con capacidad instalada en nube, integración y redes (sección 1.6).

Los otros tres proyectos aportan capacidades que no dependen del rubro. NexoPuerto resuelve el registro sin cobertura y la reconstrucción del recorrido de cada unidad. Planta Conecta lee la red de control por OPC UA sin escribir en ella. RedFrío 360 opera muchos sitios con enlaces desiguales, como los seis puntos operacionales de Aranda.

Las certificaciones institucionales cubren calidad, seguridad de la información, continuidad, servicios de TI, nube, seguridad y salud ocupacional y gestión ambiental. La Tabla \ref{tab:certificaciones-institucionales} las presenta con su organismo certificador, su vigencia y el frente del servicio en Aranda que respaldan.

\begin{tabla}{I{0.148}I{0.216}I{0.184}I{0.248}I{0.204}}{Certificaciones institucionales.}{tab:certificaciones-institucionales}
  {\cab{Norma} & \cab{Ámbito} & \cab{Organismo certificador} & \cab{Vigencia} & \cab{Frente que respalda en Aranda}}
  {Fuente: elaboración propia.}
ISO 9001:2015 & Gestión de la calidad & Bureau Veritas Chile & Hasta mayo de 2028 & Control de calidad de entregables y servicios \\
ISO/IEC 27001:2022 & Seguridad de la información & BSI Group & Hasta agosto de 2027 & Seguridad de la información de SITCO Aranda \\
ISO 22301:2019 & Continuidad del negocio & SGS Chile & Hasta noviembre de 2027 & Continuidad de la operación 24x7 \\
ISO/IEC 20000-1:2018 & Gestión de servicios de TI & SGS Chile & Hasta febrero de 2028 & Mesa de ayuda, NOC y SOC \\
ISO/IEC 27017 e ISO/IEC 27018 & Seguridad y datos personales en la nube & BSI Group, extensión de la ISO/IEC 27001 & Hasta agosto de 2027 & Servicios y datos personales en Azure \\
ISO 45001:2018 e ISO 14001:2015 & Seguridad y salud ocupacional, gestión ambiental & Bureau Veritas Chile & Hasta abril de 2028 & Trabajo del personal en la faena \\
\end{tabla}
Cada certificación respalda un frente del servicio en Aranda, desde la seguridad de la información hasta el trabajo del personal en la faena. Las certificaciones de Overfore que vencen en cualquier mes del contrato se renuevan antes de su vencimiento, con cargo a Overfore y fuera de las horas del plan de trabajo, y sus copias acompañan los antecedentes del Sobre N° 1.

Además, Overfore inició con BSI Group la certificación ISO/IEC 42001:2023 de gestión de inteligencia artificial. Su alcance son los dos componentes de inteligencia artificial de la solución: el clasificador asistido de detenciones, que opera desde la Etapa 1, y la verificación facial de portería, que entra con la Etapa 2. El plan, firmado por el representante legal, tiene tres hitos: diagnóstico de brechas en el mes 3, auditoría interna en el mes 9 y expediente de certificación en el mes 12 del contrato. Hasta obtener el certificado, esta norma no se presenta como certificación vigente.

Las certificaciones y competencias de cada integrante del equipo clave se detallan en el Anexo E del Subdocumento 1.

\seccion{Estructura para Proyecto}
Para Aranda, la empresa organiza la dotación en cuatro niveles: equipo clave, dotación de apoyo, especialistas externos y socios tecnológicos. La Figura \ref{fig:organizacion-proyecto-sitco-aranda} muestra el equipo clave y lo que conduce cada líder.

\figura{14.81}{fig-1-2-organizacion-proyecto-sitco-aranda}{Organización del proyecto SITCO Aranda.}{fig:organizacion-proyecto-sitco-aranda}{Fuente: elaboración propia.}
El Jefe de Proyecto responde ante el Comité Ejecutivo y la Contraparte Técnica del cliente. En el equipo clave, nueve personas ejercen once roles, que cubren la arquitectura, los requerimientos, los datos, la integración, el desarrollo, la seguridad, la calidad, la implantación, la operación y la gerencia del servicio. Benjamín Velásquez, además de liderar el trabajo de datos, se hace cargo de la implantación y la gestión del cambio desde el segundo mes del contrato, con la dedicación que exige la sección 19.2 de las Bases técnicas transversales (CMAA, 2026c) desde el mes 8. Claudio Cabello dirige la implementación y asume la gerencia del servicio cuando comienza la operación, en el mes 21.

Cada líder conduce a la dotación de apoyo y a los externos de su frente. Ocho desarrolladores en los meses 3 a 11 y siete en los meses 13 a 17 se reparten en los tres equipos Scrum del Subdocumento 6 (6.2.1). En el mes 12 y en los meses 18 a 20 hay dos desarrolladores, cada uno conducido por un líder del equipo clave, el equipo de plataforma, por ejemplo, responde al Líder de Integración.

El Jefe de Proyecto tiene dedicación exclusiva durante la implementación, y el Líder de Desarrollo y el Líder Funcional, dedicación completa. Los demás líderes tienen presencia permanente en su fase, y el Gerente de Servicio está dedicado a la operación. El Anexo E del Subdocumento 1 declara la dedicación y el período de cada integrante.

La Tabla \ref{tab:niveles-dotacion-proyecto} resume los cuatro niveles, su composición y dónde está su detalle.

\begin{tabla}{I{0.211}I{0.434}I{0.171}I{0.184}}{Niveles de la dotación del proyecto.}{tab:niveles-dotacion-proyecto}
  {\cab{Nivel} & \cab{Composición} & \cab{Período} & \cab{Detalle}}
  {Fuente: elaboración propia.}
Equipo clave & 9 personas que ejercen 11 roles & Según el rol, M1 a M56 & Figura \ref{fig:organizacion-proyecto-sitco-aranda} y Anexo E del Subdocumento 1 \\
Apoyo en la implementación & Hasta 30 personas: desarrollo, pruebas, datos, digitalización, terreno y mesa de ayuda & M1 a M20 & Anexo F del Subdocumento 1 \\
Apoyo en la operación & 21 personas: mesa y NOC, SOC, terreno y mantención de 26 entre los meses 21 y 23 & M21 a M56 & Anexo F del Subdocumento 1 \\
Especialistas externos & 12 perfiles para trabajos puntuales & Según el trabajo & Anexo F del Subdocumento 1 \\
Socios tecnológicos & Microsoft, NubeAndes, Integra OT y RedNorte Industrial SpA & Todo el contrato & Sección 1.6 \\
\end{tabla}
El trabajo continuo, como el desarrollo y la operación 24x7, queda en personal propio. Los especialistas externos se contratan solo para trabajos puntuales, como las pruebas de intrusión, el estudio de cobertura o el montaje de la sala técnica.

Cada contratación externa requiere autorización previa y escrita del cliente, y el conjunto de las subcontrataciones se mantiene bajo el 40 \% del valor del contrato.

Todo el personal que sube a la faena, sea del equipo clave, de apoyo o externo, se acredita como contratista antes de su primer turno, con exámenes preocupacionales, examen de altura geográfica, inducción y credencial. La faena está a 3.180 m de altura y la acreditación toma alrededor de nueve días (CMAA, 2026b, cap. 4, p. 9, y cap. 15, p. 26). Por eso los técnicos de terreno y los digitalizadores entran acreditados al mes en que comienza su trabajo.

\seccion{Alianzas}
La empresa mantiene cuatro alianzas tecnológicas vigentes. En todas, responde ante el cliente por el servicio completo, y cada socio trabaja dentro de un alcance declarado. La Tabla \ref{tab:alianzas-tecnologicas-vigentes} presenta con el alcance general de cada alianza.

\begin{tabla}{I{0.185}I{0.233}I{0.340}I{0.243}}{Alianzas tecnológicas vigentes.}{tab:alianzas-tecnologicas-vigentes}
  {\cab{Socio} & \cab{Tipo} & \cab{Alcance de la alianza} & \cab{Evidencia}}
  {Fuente: elaboración propia.}
Microsoft Chile Ltda. (Microsoft Azure) & Microsoft Solutions Partner for Infrastructure (Azure) & Plataforma de nube pública: diseño, despliegue y soporte & Certificado del proveedor y carta de compromiso \\
NubeAndes Cloud SpA & Socio de servicios administrados sobre Microsoft Azure & Administración de los servicios desplegados en Azure y apoyo a la continuidad operacional, bajo responsabilidad de Overfore & Carta de compromiso con alcance de administración y continuidad \\
Integra OT Ltda. & Socio TI/TO & Integración industrial y protocolos OPC UA, sin modificar PLC ni DCS & Carta de colaboración \\
RedNorte Industrial SpA & Socio de redes y terreno & Redes industriales, salas técnicas, gabinetes y estudios de cobertura & Carta de colaboración \\
\end{tabla}
Microsoft Azure es la plataforma de nube pública de SITCO Aranda. NubeAndes aporta servicios de administración sobre esa plataforma y apoyo a la continuidad operacional. Integra OT cubre la integración industrial y RedNorte, las redes y el trabajo en terreno. Overfore mantiene la responsabilidad por el servicio completo y acredita su condición de Solutions Partner de Microsoft mediante el certificado indicado en la Tabla \ref{tab:alianzas-tecnologicas-vigentes}.

Los certificados y las cartas de los cuatro socios se acompañan en el Sobre N° 1, con los antecedentes técnicos habilitantes. En Aranda, RedNorte Industrial SpA participa además como subcontratista, con el detalle en el Anexo F del Subdocumento 1; su régimen de control se desarrollará en el Subdocumento 12 de la Propuesta Final.

\seccionsinnumero{Referencias}
\begin{referencias}
\item Compañía Minera Altos de Aranda. (2026b). \emph{Bases técnicas del caso Minería}. Licitación N° TFEP-01/2026.
\item Compañía Minera Altos de Aranda. (2026c). \emph{Bases técnicas transversales}. Licitación N° TFEP-01/2026.
\end{referencias}

\seccionsinnumero{Declaración de uso de IA}
Se utilizó IA generativa, principalmente como apoyo en la redacción y en la organización de tablas y organigramas; el detalle por sección está en la Tabla \ref{tab:declaracion-uso-ia-seccion}.

\begin{tabla}{I{0.175}I{0.163}I{0.190}I{0.088}I{0.141}I{0.245}}{Declaración de uso de IA por sección.}{tab:declaracion-uso-ia-seccion}
  {\cab{Sección} & \cab{Herramienta} & \cab{Finalidad del uso} & \cab{Nivel en\newline texto} & \cab{Nivel en\newline diagramas} & \cab{Revisión humana (quién y qué verificó)}}
  {Fuente: elaboración propia.}
Introducción y capítulo 1 & Claude & Borrador de presentación de Overfore y del propósito del capítulo. & Alto & Ninguno & Kavon Kermani: validó el alcance de la presentación y su coherencia con la oferta. \\
1.1 & claude & Apoyo para ordenar la descripción de la empresa y sus servicios. & Medio & Ninguno & Benjamín Rodríguez: revisó que los antecedentes de la empresa coincidan con el resto de la propuesta. \\
1.2 & claude & Organización del texto que acompaña la estructura de la empresa. & Medio & Medio & Kavon Kermani: verificó los cargos, dependencias y nombres del organigrama. \\
1.3 & claude & Reformulación de las prácticas internas de calidad, seguridad y conocimiento. & Medio & Ninguno & Kavon Kermani: aseguró que las prácticas descritas tengan respaldo en el capítulo. \\
1.4 & claude & Síntesis de la experiencia y de las certificaciones mencionadas. & Medio & Ninguno & Marcelo Flores: revisó las referencias de proyectos y el alcance de cada certificación. \\
1.5 & claude y chatGPT & Base de redacción para explicar la organización del proyecto. & Medio & Medio & Kavon Kermani: comprobó la dotación y las responsabilidades frente a T-14 y T-15. \\
1.6 & claude & Ajustes de claridad en la descripción de las alianzas. & Medio & Medio & Benjamín Rodríguez: revisó los servicios aportados por cada aliado y sus límites en el proyecto. \\
Referencias & claude & Apoyo para ordenar y uniformar las referencias citadas. & Medio & Ninguno & Marcelo Flores: validó la existencia de las fuentes y su correspondencia con las citas. \\
Anexo A del Subdocumento 1 & claude & Ajustes puntuales al texto de la carta de Terminal Logístico Bahía Sur. & Bajo & Ninguno & Marcelo Flores: comprobó los antecedentes, fechas y alcance de la referencia. \\
Anexo B del Subdocumento 1 & claude & Apoyo de claridad en la carta de Alimentos Cordillera. & Bajo & Ninguno & Marcelo Flores: revisó la correspondencia entre la carta y la experiencia declarada. \\
Anexo C del Subdocumento 1 & claude & Revisión de redacción de la carta de Frío Austral. & Bajo & Ninguno & Marcelo Flores: analizó los datos del proyecto y de la entidad de referencia. \\
Anexo D del Subdocumento 1 & claude & Ajustes de forma al borrador de la carta de Minera Pampa Tarel. & Bajo & Ninguno & Benjamín Rodríguez: revisó que los antecedentes de la carta coincidan con el T-6. \\
Anexo E del Subdocumento 1 & claude & Organización de los perfiles del equipo clave. & Medio & Ninguno & Benjamín Rodríguez: verificó que cada rol del equipo clave tenga la experiencia y la certificación que exige la sección 19.2 de las Bases técnicas transversales. \\
Anexo F del Subdocumento 1 & claude & Apoyo para describir la dotación de apoyo y los especialistas. & Medio & Ninguno & Benjamín Rodríguez: verificó la función y la dedicación de cada perfil de apoyo frente al T-15. \\
T-6 & claude & Ajustes puntuales a la presentación de experiencia en proyectos similares. & Bajo & Ninguno & Kavon Kermani: revisó los proyectos, referencias y antecedentes frente a los anexos. \\
\end{tabla}

\end{document}
